\maketitle
\chapter{Il pacchetto Python}

Come richiesto dalla traccia, l'intero sistema è reso disponibile come pacchetto Python nativo, denominato \texttt{Gruppo\_Ferrari\_DeFusco\_Cuconato}. Il requisito fondante della progettazione è stato \textbf{non duplicare l'algoritmo}: il pacchetto non reimplementa nulla, ma espone attraverso le Python C-API lo stesso identico nucleo di calcolo C validato nei capitoli precedenti (nella sua variante a \emph{intrinseci SIMD}, portabile su ogni piattaforma). Ne consegue che ogni proprietà dimostrata per gli eseguibili — correttezza rispetto ai riferimenti ufficiali inclusa — vale automaticamente anche per la libreria.

\section{Interfaccia e moduli}
Il pacchetto espone tre sotto-moduli, uno per configurazione hardware, ciascuno con la medesima classe \texttt{QuantPivot} dotata di un'interfaccia in stile \emph{scikit-learn}:

\begin{table}[h!]
\centering
\small
\renewcommand{\arraystretch}{1.4}
\begin{tabularx}{\textwidth}{|l|l|X|}
\hline
\textbf{Modulo} & \textbf{Precisione} & \textbf{Back-end} \\ \hline
\texttt{quantpivot32}    & \texttt{float32} & SIMD SSE2 (intrinseci) \\ \hline
\texttt{quantpivot64}    & \texttt{float64} & SIMD AVX2 (intrinseci) \\ \hline
\texttt{quantpivot64omp} & \texttt{float64} & SIMD AVX2 (intrinseci) + OpenMP \\ \hline
\end{tabularx}
\caption{I tre moduli del pacchetto.}
\end{table}

Il metodo \texttt{fit(dataset, n\_pivots, quant\_level)} costruisce l'indice a pivot (quantizzazione del dataset e matrice $\tilde{d}(v,p)$) e restituisce l'oggetto stesso, consentendo il concatenamento; \texttt{predict(query, k)} esegue la ricerca dei $k$ vicini per tutte le query e restituisce la tupla \texttt{(ids, dists)}: due array NumPy di forma $(n_q, k)$ contenenti gli identificativi dei vicini e le rispettive distanze euclidee reali.

\begin{verbatim}
import numpy as np
from Gruppo_Ferrari_DeFusco_Cuconato.quantpivot64omp import QuantPivot

DS = np.ascontiguousarray(dataset, dtype=np.float64)   # (N, D)
Q  = np.ascontiguousarray(query,   dtype=np.float64)   # (nq, D)

model      = QuantPivot().fit(DS, n_pivots=16, quant_level=64)
ids, dists = model.predict(Q, k=8)
\end{verbatim}

\section{Architettura a tre strati}
L'integrazione è organizzata in tre strati, dal più esterno al più interno:

\begin{enumerate}
    \item \textbf{Wrapper C-API} (\texttt{quantpivot*\_py.c}): definisce il tipo Python \texttt{QuantPivot} e i suoi metodi. Riceve gli array NumPy, ne verifica forma (bidimensionale), tipo (\texttt{float32}/\texttt{float64} a seconda del modulo) e contiguità in memoria, quindi ne passa al livello inferiore i puntatori ai dati grezzi: \textbf{nessuna copia} dei dati viene mai effettuata.
    \item \textbf{Adattatore} (\texttt{quantpivot32.c}, \texttt{quantpivot64.c}, \texttt{quantpivot64omp.c}): traduce la struttura ponte \texttt{params} nelle strutture native del nucleo (\texttt{MatrixF32}/\texttt{MatrixF64}, \texttt{Index}, \texttt{Neighbor}) e invoca \texttt{build\_index()} in \texttt{fit} e \texttt{knn\_query\_all()} in \texttt{predict}.
    \item \textbf{Nucleo condiviso}: gli stessi \texttt{index.c}, \texttt{query.c}/\texttt{query64.c}, \texttt{quantization.c} e il calcolo della distanza vettorizzato con gli \emph{intrinseci SIMD} di \texttt{distance.c}, compilato con le macro \texttt{USE\_SSE2} (32~bit) o \texttt{USE\_AVX} (64~bit), più \texttt{-fopenmp} per la variante parallela.
\end{enumerate}

La struttura \texttt{params}, definita in \texttt{include/common.h}, è l'unico oggetto che attraversa il confine tra gli strati: raccoglie i puntatori a dataset e query, i parametri del problema ($h$, $k$, $x$, dimensioni), il puntatore all'indice costruito e i buffer dei risultati. La macro \texttt{QP\_DOUBLE} commuta il tipo scalare \texttt{type} (\texttt{float}/\texttt{double}) e la costante di allineamento \texttt{align} (16/32 byte), permettendo di condividere lo stesso wrapper tra le versioni a diversa precisione.

\section{Gestione della memoria tra i due mondi}
La coesistenza del garbage collector di Python con le allocazioni manuali del C richiede alcune cautele, tutte gestite nel wrapper:

\begin{itemize}
    \item \textbf{Riferimenti agli input}: \texttt{fit} e \texttt{predict} incrementano il contatore di riferimenti (\texttt{Py\_INCREF}) degli array NumPy ricevuti e lo rilasciano solo alla distruzione dell'oggetto: gli array non possono essere deallocati da Python mentre il C ne sta usando i puntatori.
    \item \textbf{Proprietà dei risultati}: i buffer di \texttt{ids} e \texttt{dists} sono allocati in C (\texttt{\_mm\_malloc}) e ``adottati'' dagli array NumPy restituiti tramite \texttt{PyCapsule} con distruttore: quando l'ultimo riferimento Python all'array cade, il distruttore libera la memoria C. Non esistono quindi né copie né perdite di memoria.
    \item \textbf{Distruzione del modello}: alla deallocazione dell'oggetto \texttt{QuantPivot}, l'indice viene rilasciato con la stessa \texttt{free\_index()} usata dagli eseguibili.
\end{itemize}

\section{Sistema di build}
Il pacchetto si compila con \texttt{setuptools} (\texttt{setup.py} + \texttt{pyproject.toml} nella radice del progetto), che dichiara tre estensioni native — una per modulo — composte dal wrapper, dall'adattatore e dal nucleo di calcolo condiviso. Per garantire la \textbf{portabilità}, il calcolo della distanza passa per gli \emph{intrinseci SIMD} (\texttt{distance.c}) anziché per le routine Assembly esterne: gli intrinseci sono supportati in modo identico da \texttt{gcc}, \texttt{clang} e MSVC, quindi il pacchetto si compila con il \emph{compilatore di sistema} su Linux, macOS e Windows, senza richiedere alcun assemblatore né toolchain specifico.

Una classe \texttt{build\_ext} personalizzata seleziona automaticamente i flag corretti in base al compilatore rilevato: \texttt{-O3 -msse2}/\texttt{-mavx2} e \texttt{-fopenmp} per \texttt{gcc}/\texttt{clang}, oppure \texttt{/O2 /arch:AVX2 /openmp} per MSVC. Nell'ambiente di riferimento della traccia (Linux, \texttt{gcc}) l'installazione si riduce a \texttt{pip install .} dalla cartella del progetto: il compilatore genera le estensioni e collega OpenMP tramite la \texttt{libgomp} di sistema, senza passi aggiuntivi.

Il posizionamento di \texttt{setup.py} nella radice (con il pacchetto sotto \texttt{python/}) evita inoltre il problema di \emph{shadowing}: dopo l'installazione l'\texttt{import} risolve sui moduli compilati in \texttt{site-packages} e non sui sorgenti, privi delle estensioni.

Le routine Assembly scritte a mano (\texttt{distance\_sse2.S}, \texttt{distance\_avx2.S}) restano parte integrante del progetto: alimentano gli eseguibili standalone e il confronto prestazionale del Capitolo~4.

\section{Verifica di correttezza del pacchetto}
La libreria è stata sottoposta alla stessa verifica degli eseguibili: per ciascuno dei tre moduli, \texttt{fit} + \texttt{predict} sul carico di lavoro di riferimento ($n=2000$, $D=256$, $h=16$, $k=8$, $x=64$) riproducono \textbf{2000 query su 2000} dei file di riferimento ufficiali — stessi identificativi nello stesso ordine e stesse distanze (tolleranza $10^{-3}$ per \texttt{float32}, $10^{-9}$ per \texttt{float64}). Il percorso verificato attraversa l'intera catena: validazione NumPy, wrapper C-API, adattatore e nucleo C con calcolo vettorizzato tramite intrinseci SIMD (con OpenMP attivo nella terza variante).

Lo script dimostrativo \texttt{examples/esempio\_knn.py} ripete questa verifica a ogni esecuzione, riportando anche i tempi: sul carico di riferimento la variante \texttt{quantpivot64omp} completa la fase di \texttt{predict} in circa 60--70~ms, coerentemente con i risultati del Capitolo 4. Il beneficio del multi-threading resta dunque pienamente fruibile anche dall'interno dell'interprete Python: durante l'esecuzione delle routine native il GIL non costituisce un ostacolo, poiché il parallelismo è gestito interamente da OpenMP nello strato C.
